\chapter{Provisioning}
\label{ch:provisioning}

\section{Block devices}
\label{sec:block-devices}

Each hard disk or hardware RAID set appears as a so-called block device.

\begin{itemize}
\item
  UNIX-like systems show block devices in the \texttt{/dev} directory:

  \begin{itemize}
  \item
    Linux show SATA devices as \texttt{sd} and a letter, such as
    \texttt{/dev/sda}, \texttt{/dev/sdb}.
  \item
    Mac systems show \texttt{disk} followed by a number, like
    \texttt{/dev/disk0}.
  \item
    Varies on other unix-like systems (e.g.~Xenix, AIX, BSD)
  \end{itemize}
\item
  Windows systems will show block devices in the logical disk manager.
  They will \emph{NOT} appear as drive letters!
\end{itemize}

A filesystem is then setup on the block device by formatting it. The
filesystem organises the block device into the familiar structure of
files and folders. Also the filesystem typically handles permissions and
metadata storage.

\section{Partitioning}
\label{sec:partitioning}

Partitioning allows us to divide a single block device into a number of
separate segments that each appear as a separate block device. Each
partition can have a different filesystem on it. In almost all cases
fixed disks in PC-based systems are partitioned.

There are a number of partitioning schemes. A few common concepts:

\begin{itemize}
\item
  A partition table defines one or more partitions on the disk. This is
  located in a known location, usually first block on drive.

  \begin{itemize}
  
  \item
    \href{https://en.wikipedia.org/wiki/Partition_type}{Partition
    types}.
  \end{itemize}
\item
  The host's OS needs to boot usually from a disk. Partition scheme
  needs to be compatible with the BIOS or UEFI to achieve this. Also the
  boot loader may need to be set up.
\end{itemize}

Most PC-based systems use either MBR or GPT.

\subsection{Master Boot Record (MBR)}
\label{sec:master-boot-record-mbr}

MBR is the most common PC partitioning scheme. Introduce circa 1983 with
DOS. This defines up to four primary partitions on a drive, .

\autoimage{mbr_scheme}{MBR partition scheme}{mbr-scheme}

\begin{itemize}
\item
  MBR layout closely linked with BIOS boot process. One partition marked
  ``active''.
\item
  The MBR scheme stores the partition table at the beginning of the
  first 512 bytes of the drive.

  \begin{itemize}
  
  \item
    446 bytes bootloader, 64 bytes partition table, 2 bytes ``magic
    number''
  \end{itemize}
\item
  MBR limited to drives less than approx 2.2 TB.
\end{itemize}

\subsection{Extended Boot Record (EBR)}
\label{sec:extended-boot-record}

EBR allows more than four partitions on a drive. One MBR partition is
designated as an extended partition. This \emph{extended partition} can
hold one or more \emph{logical partitions}, .

\autoimage{ebr_scheme}{EBR scheme}{ebr-scheme}

\subsection{GUID Partition Table (GPT)}
\label{sec:guid-partition-table}

GPT is a more modern scheme and is associated with UEFI-based systems.
UEFI is an alternative to BIOS and is used by newer PCs and Apple Macs.
GPT layout is shown in .

\autoimage{gpt_scheme}{GPT partition table scheme}{gpt-scheme}

Points to note:

\begin{itemize}
\item
  GPT can define up to 128 partitions.
\item
  A so-called \emph{protective MBR} is present in the usual MBR
  location:

  \begin{itemize}
  \item
    Indicates entire disk is a single MBR partition.
  \item
    Shows disk partitioning (and other) software that the MBR scheme is
    not in use.
  \item
    Prevents accidental overwriting by MBR utilities.
  \end{itemize}
\item
  GPT partition table follows protective MBR.
\item
  GPT stores redundant copy of partition table at the end of the drive.
\item
  Boot process more complicated than MBR.
\end{itemize}

\section{Logical volumes}
\label{sec:logical-volumes}

Logical Volumes can sometimes offer a practical alternative to
partitions. Partitioning can be considered as a thick provisioning of
disk space. Logical volumes represent a thin provisioning of disk space.

Unlike partitoning, logical volumes are OS-dependent and their handling
depends on the host OS in use. The general ideas are:

\begin{itemize}
\item
  Volumes are created which appear to OS as a partition, which are then
  formatted and mounted in the usual way.
\item
  A single disk is used in its entirety. Usually a single large
  partition is created to hold the logical volumes. Sometimes a small
  standard boot partition is required.
\item
  Logical Volume manager can resize volumes. For this to work
  non-destructively, file system must support resizing. Useful
  capability when used appropriately.
\item
  May support snapshot/restore operations.
\end{itemize}

\subsection{Linux LVM}
\label{sec:linux-lvm}

Linux has an inbuilt Logical Volume Manager (LVM). LVM is built into the
kernel and managed by a number of commands in \texttt{/usr/sbin}.

LVM takes over partitions on physical disks as physical volumes.
Physical volumes are aggregated into volume groups. Logical Volumes,
which act like a partition, are then provisioned on a volume group.

\begin{description}
\item[Physical disks]
appear to the system as block devices.
\item[LVM partitions]
are created on the physical disk, usually 1 for entire disk. 

\item[Physical volumes]
map to physical disks. Physical volumes are created on a partition using
\texttt{pvcreate}.

Commands dealing with physical volumes start with \texttt{pv}. For
example, we can use \texttt{pvdisplay} to show information about all
physical volumes in the system.

\item[Volume group]
aggregates one or more physical volumes to create a virtual hard disk
using \texttt{vgcreate}. The volume group has a name / label, here
\texttt{vg-01}.

Commands dealing with volume groups start with \texttt{vg}. For example,
we can use \texttt{vgdisplay} to show information about volume groups in
the system. The volume group will then appear as a directory under
\texttt{/dev}, in this case \texttt{/dev/vg01}.

\item[Logical volume]
is provisioned on the volume group, and acts as a virtual hard disk.

Commands dealing with logical volumes start with \texttt{lv}. For
example, we can use \texttt{lvdisplay} to show information about logical
volumes in the system. The new logical volume will then appear in the
\texttt{/dev/vg01} folder as \texttt{/dev/vg01/vol\_data}.
\end{description}

\autoimage{lvm}{LVM}{lvm}

Example of
creating an LVM partition on entire of \texttt{/dev/sda}:

\inputminted{bash}{lvm_commands.sh}

\section{Useful links}\label{useful-links}

\begin{itemize}
\item
  \url{https://www.digitalocean.com/community/tutorials/an-introduction-to-lvm-concepts-terminology-and-operations}
\item
  \url{https://www.server-world.info/en/note?os=Ubuntu_17.04\&p=iscsi\&f=1}
\item
  \url{https://en.wikipedia.org/wiki/Logical_volume_management}
\end{itemize}
